\section{\label{sec:influence_to_hmf}From Influence Functionals to the Hamiltonian of Mean Force}

We now eliminate the harmonic bath exactly and express the reduced equilibrium operator as a Euclidean path integral. For Gaussian baths with linear coupling, this elimination is exact and produces an influence functional that encodes all bath effects in a single memory kernel. The object of interest is the reduced equilibrium kernel in the system position basis,
\begin{equation}
    \bar\rho_Q(q_f,q_i;\beta) := \langle q_f|\Tr_X e^{-\beta H_{\mathrm{tot}}}|q_i\rangle.
    \label{eq:reduced_kernel_again}
\end{equation}
Following the influence-functional formalism \cite{feynmanTheoryGeneralQuantum1963a,caldeiraQuantumTunnellingDissipative1983a,grabertQuantumBrownianMotion1988} and specifically its stochastic representation for Gaussian baths \cite{mccaulPartitionfreeApproachOpen2017c, mccaulDrivingSpinbosonModels2018a, mccaulHowWinFriends2021b}, the resulting reduced density matrix admits the imaginary-time path integral representation:
\begin{widetext}
\begin{equation}
    \bar\rho_Q(q_f,q_i;\beta) = Z_X(\beta) \int_{q(0)=q_i}^{q(\beta)=q_f}\!\!\mathcal D q\; \exp\Bigg[ -S_Q[q] + \frac{1}{2}\int_0^\beta d\tau\int_0^\beta d\tau' f\big(q(\tau)\big) K(\tau-\tau') f\big(q(\tau')\big) \Bigg].
    \label{eq:reduced_effective_PI}
\end{equation}
\end{widetext}
Here $Z_X(\beta)=\Tr_X e^{-\beta H_X}$ is the free bath partition function, $S_Q$ is the Euclidean system action defined in Sec.~\ref{sec:quenched}, and $K(\tau)$ is the imaginary-time influence kernel. In terms of the bath spectral density $J(\omega)$, the kernel is given by
\begin{equation}
    K(\tau) = \frac{1}{\pi}\int_0^\infty d\omega\, J(\omega) \frac{\cosh\!\left[\omega\left(\frac{\beta}{2}-|\tau|\right)\right]} {\sinh\!\left(\frac{\beta\omega}{2}\right)}.
    \label{eq:kernel_spectral}
\end{equation}

While the representation \eqref{eq:reduced_effective_PI} is exact, it is nonlocal (bilocal) in imaginary time. This nonlocality can be resolved by a Hubbard--Stratonovich (HS) transformation \cite{hubbardCalculationPartitionFunctions1959a, stratonovich1957QDistro}, which decouples the bilocal interaction by introducing an auxiliary Gaussian field $\xi(\tau)$. This step provides the bridge to an operator-level representation of $\bar{\rho}_Q$ as a Gaussian average over the local quenched canonical densities discussed in Sec.~\ref{sec:quenched}:
\begin{equation}
\begin{split}
    \bar\rho_Q(\beta) 
    &= 
    Z_X(\beta)\,\mathbb E_\xi\bigg[ \\
    &\quad \hat{\mathcal{T}}\exp\!\left(-\int_0^\beta d\tau\,\bar{H}(\tau;\xi)\right) \bigg],
\end{split}
    \label{eq:bar_rho_as_quenched_average}
\end{equation}
where the stochastic process $\xi(\tau)$ is a zero-mean Gaussian field with covariance determined by the influence kernel \cite{mccaulPartitionfreeApproachOpen2017c,mccaulDrivingSpinbosonModels2018a,mccaulHowWinFriends2021b}
\begin{equation}
    \mathbb{E}_\xi[\xi(\tau)] = 0, \quad \mathbb{E}_\xi[\xi(\tau)\xi(\tau')] = K(\tau-\tau').
    \label{eq:noise_statistics}
\end{equation}
This is the promised bridge: the nonlocal influence-functional path integral is exactly equivalent to an average over quenched, system-local imaginary-time propagators.

In this form, finding the $H_{\mathrm{MF}}(\beta)$ is reduced to performing the Gaussian integration over $\xi$. Notably, the barrier to evaluating this does \emph{not} lie in the usual difficulties associated with calculating $H_{\mathrm{MF}}(\beta)$, where it is strong coupling/low temperature that are identified as the main obstacle. Instead, the difficulty lies in the algebraic structure of $H_Q$ and $f$. This is the key insight that allows us to find the HMF exactly in the commuting sector. Imposing
\begin{equation}
    [H_Q,f]=0,
    \label{eq:commuting_here}
\end{equation}
the generator $\bar{H}(\tau;\xi) = H_Q - \xi(\tau) f$ (the sign is conventional; since $\xi$ is Gaussian with zero mean, $\xi\to-\xi$ leaves the average unchanged) now commutes with itself at different imaginary times. The $\hat{\mathcal{T}}$ ordering becomes redundant and the expression factorises:
\begin{align}
    \hat{\mathcal{T}}\exp\!\left[-\int_0^\beta d\tau\,\bar{H}(\tau;\xi)\right]
    &=
    e^{-\beta H_Q}\, \exp\!\left[f\,X\right],
    \label{eq:tilde_rho_commuting_factorised}
\end{align}
where $X:=\int_0^\beta d\tau\,\xi(\tau)$ is a Gaussian variable with mean zero and variance
\begin{equation}
    C(\beta) = \mathbb E_\xi[X^2] = \int_0^\beta d\tau\int_0^\beta d\tau'\,K(\tau-\tau').
    \label{eq:Cbeta_def}
\end{equation}
Substituting the kernel \eqref{eq:kernel_spectral} and performing the double
imaginary-time integral gives
\begin{align}
    C(\beta)
    &=\frac{2\beta}{\pi}\int_0^\infty d\omega\,\frac{J(\omega)}{\omega}
    \equiv 2\beta\,\lambda,
    \nonumber\\
    \lambda&:=\frac{1}{\pi}\int_0^\infty d\omega\,\frac{J(\omega)}{\omega},
    \label{eq:Cbeta_explicit}
\end{align}
where $\lambda$ is the reorganisation energy.
Crucially, $C(\beta)$ is \emph{linear} in $\beta$, so the ratio
$C(\beta)/(2\beta)=\lambda$ is temperature independent.

Therefore the Gaussian average $\mathbb E_\xi[e^{f X}] = \exp[\frac{1}{2}C(\beta)f^2]$ can be performed exactly, and hence, using \eqref{eq:bar_rho_as_quenched_average},
\begin{equation}
    \bar\rho_Q(\beta) = Z_X(\beta)\,e^{-\beta H_Q}\,\exp\!\left[\beta\lambda\,f^2\right].
    \label{eq:bar_rho_commuting_closed}
\end{equation}
The Hamiltonian of mean force follows by taking
$-\frac{1}{\beta}\log\bar\rho_Q$:
\begin{equation}
    H_{\mathrm{MF}}(\beta) = H_Q -\lambda\,f^2 +\mathrm{const}(\beta)\,\mathbf 1,
    \label{eq:HMF_commuting_final}
\end{equation}
where $\mathrm{const}(\beta)$ is a scalar that is fixed by the normalisation
$\Tr\,e^{-\beta H_{\mathrm{MF}}}=1$ and therefore drops out of the
normalised mean-force Gibbs state
$\rho_{\mathrm{MF}}=e^{-\beta H_{\mathrm{MF}}}/\Tr\,e^{-\beta H_{\mathrm{MF}}}$.
The only nontrivial operator content is the correction $-\lambda f^2$, which
is temperature independent.

Note that $\lambda$ coincides with the static potential renormalisation
$\frac{1}{2}\sum_k c_k^2/(m_k\omega_k^2)$ of the Caldeira--Leggett
model. If the standard counterterm $+\lambda f^2$ has been absorbed into
$H_Q$ (as stated in Sec.~\ref{sec:model}), the bath-induced shift
$-\lambda f^2$ cancels it exactly:
$H_{\mathrm{MF}}$ reduces to the bare system Hamiltonian $H_Q^{\mathrm{bare}}$,
and the normalised mean-force Gibbs state coincides with the bare-system
Gibbs state. In other words, in the commuting sector with the CL counterterm
convention, the bath has no operator-level effect on the reduced equilibrium
state---its only role is to shift the total free energy.
Conversely, if no counterterm is included, $-\lambda f^2$ represents a genuine
bath-induced static renormalisation of the system potential.
